\Needspace{12\baselineskip}
\originalchapter{十九份实验手册的采集与分析链}
前文建立了每个实验的物理模型。此章将实际公开手册的测量要求接到原始数据、校准、拟合与判据，供自学者设计报告。所有现场采集在本项目均未执行；以下是有条件的分析方案，未知的仪器参数和实际观测结果保留为空缺。原手册的讨论/口试主题与正式预习题分别记录，不将每条操作步骤凑作一道考试题。

\originalsection{expI：光电效应的截止定义}
汞灯稳定后用各滤片照射光阴极，记录完整 $I(V)$，越过零流继续取点以识别反向电流；手册要求至少五次重复扫描。光若照到阳极环会产生反向电流，环境光、接地和灯漂移也会改变尾部。原始表应包括滤片、偏压、流、重复号与暗流，而不只剩一列截止电压。

用暗流扣除后的固定小流准则和带背景的曲线外推两种方法分别确定截止，比较得到h的变化作为模型系统误差线索。滤片中心频率及带宽进入横轴，校准误差由各线共享；由 $V_s=(h/e)\nu-\phi/e$ 拟合斜率、截距。强度变化主要改变流，并非必然改变光子最高能量。口试中的经典蓄能估算可用原条件强度 $10^{-3}$W/m$^2$、原子截面积尺度 $(.3\mathrm{nm})^2$、逸出功4eV，得时间约 $7.1\times10^3$s；快速响应反对这种逐原子缓慢攒能图景，但单一估算并非对所有经典电磁作用的通用反证。

\originalsection{expII：四档计数与模拟分布}
手册要求平均约1、5、10、100计数每秒四档（目标约±25\%），相同门宽反复统计，记录顺序以检查率漂移；脉冲需经过闪烁体、光电倍增、前置和成形放大、阈值计数，先看示波器的极性、上升和恢复。改变源距离或门宽不等于改变基本事件概率律。

各档估 $\hat\lambda=\bar n$、方差、直方图，与Poisson $p_n=\e^{-\lambda}\lambda^n/n!$ 比较，低概率尾合并或直接似然。另用程序产生同样样本量和平均值的模拟作比较；模拟直方图必须独立标注。均值和方差不一致可能来自死时间、漂移、阈值或相关事件，不能先把不符合的样本删掉。口试要求的等待时间密度为 $Re^{-Rt}$；binomial到Poisson取N大、p小、Np固定，高均值时再近似Gaussian，三者条件不同。

\originalsection{expIII：Michelson的两通道电压}
两臂近等长，使两束在光电探测器重叠；示波器一通道读光电压、另一通道读同时驱动PZT的发生器电压。扫描上升与下降分别标定，取多个条纹周期来降低读点误差，检查振动和PZT迟滞。手册示例位移系数41.6±1.3nm/V不是自己装置必然的标定。

镜位移 $\Delta\ell=\kappa\Delta V$，光程变化是两倍，N条纹给 $\lambda=2\kappa\Delta V/N$。条纹对电压可拟合 $A+B\cos(4\pi\kappa V/\lambda+\phi)$；几次扫描共享κ，不把κ误差按重复数缩小。臂差超过相干长度会降可见度，错位也会降可见度，二者需用调节与控制分别判别。前文模拟读数仅示误差传播。

\originalsection{exp1：Compton能量与截面}
用已知gamma线标定多道分析器，将源、靶、散射探测器置于已记录的角度/距离；手册区分能量位移与微分截面两项目标。每个角度同时保留峰谱、背景谱、活时间和几何。峰心与 $E'=E/[1+(E/m_ec^2)(1-\cos\theta)]$ 比较，角宽与有限能量分辨共同展宽，不能把峰宽都归于电子静质量不确定。

截面计数要除入射光子通量、靶电子数、接受立体角、探测效率与活时间，还需入射/出射衰减。更完整的单散射预测是 $N=t\int\Phi(\mathbf r) n_e(\mathbf r)(\dd\sigma/\dd\Omega)\epsilon(E')\e^{-\mu_0l_0-\mu_1l_1}\dd V\dd\Omega$。Klein--Nishina给
\[
 \frac{\dd\sigma}{\dd\Omega}=\frac{r_e^2}{2}(E'/E)^2\left(E'/E+E/E'-\sin^2\theta\right).
\]
绝对归一参数没测全时只能报告相对角分布；多重散射/偶然符合控制会检验单散射前提，不能拿一条理论线替代几何积分。

\originalsection{exp7：Franck--Hertz与Ramsauer--Townsend}
Franck--Hertz扫加速电压，保留温度、收集反偏及全部电流曲线；相邻谷间距近似原子最低可激发能除e，接触势主要移整体曲线。重复温度与扫向可检验碰撞概率、空间电荷和漂移，不能把每个谷相对零电压的距离直接当激发能。用跨多个谷的编号对电压拟合比单间距稳健。

手册还包含氙散射的Ramsauer部分。参考空管流 $I_0(V)$ 与含气流 $I(V)$ 给 $\ln[I_0/I]\simeq nl\sigma$，原装置几何使并非所有散射都移出束，所得简单截面会偏低。波函数在势区的相位匹配能让低能散射振幅近消失；用部分波 $\sigma=(4\pi/k^2)\sum_l(2l+1)\sin^2\delta_l$ 说明极小可来自主导s波相移接近整数π。它与Franck--Hertz非弹性阈值的机制不同。未量得nl时保留相对曲线和极小位置，不捏造绝对截面。

\originalsection{exp9：电子动量、能量及球面磁场}
磁谱仪由曲率半径和场确定 $p=eBR$，探测器脉冲给动能K；分别检验非相对论 $p^2=2mK$ 与 $p^2c^2=K^2+2Kmc^2$。同一场校准影响全部p，孔宽与准直决定接受动量范围；探测器能量损失、背散射、磁场不均匀应进入残差分析。源谱是连续的，不能只挑符合理论的一点。

公开附录要求补球面线圈场。若球半径R、内部 $\mathbf B=B_i\hat z$，磁标势内为 $-B_ir\cos\theta$、外为 $C\cos\theta/r^2$；径向B连续给 $C=B_iR^3/2$。球面外 $B_\theta=(B_i/2)\sin\theta$，内 $B_\theta=-B_i\sin\theta$；按SI $\hat r\times(\mathbf B_{\rm out}-\mathbf B_{\rm in})=\mu_0\mathbf K$，得 $K_\phi=3B_i\sin\theta/(2\mu_0)$。这给连续理想线圈分布，有限匝与支架导致的偏差仍需场图测量。

\originalsection{exp11：磁场向量与泵浦恢复}
先校三轴线圈常数和零点，改变某轴电流时拟合 $\nu^2=[g_F\mu_B/h]^2[(B_x+\alpha_xI_x)^2+B_y^2+B_z^2]$；扫最小频率可以消掉该轴分量，但其他两轴未补偿时最小值不为0。双同位素斜率比提供不依绝对磁场的独立检查，弱场与强场超精细解耦范围分开。

手册还要求反转场或断/复光观察透射恢复。对小扰动可拟合人口方程 $\dot P=R(P_{\rm pump}-P)-\Gamma(P-P_{\rm eq})$，恢复时间 $1/(R+\Gamma)$；黑暗中R为0才测裸弛豫。改变光强分离泵浦率与碰撞率，扫速太快产生迟滞，不将仪器低通时间常数当原子弛豫。每组保留偏振、温度、三轴场、光强与时间序列。

\originalsection{exp12：脉冲NMR的三个时间尺度}
校共振频率和90/180度脉宽后，记录复FID和脉冲时序。$\pi/2-\tau-\pi$在 $2\tau$附近回波；静态场不均匀能重聚相位，因此FID给 $T_2^*$，回波振幅对 $2\tau$ 衰减给T2，不能互换。180度脉冲失准、扩散与梯度也会影响回波。

T1用反转恢复 $M_z(t)=M_0(1-2\e^{-t/T_1})$ 或饱和恢复，读出相位须保存，取绝对值会在过零点造成模型改变。按手册比较样品/水甘油浓度或黏度时，把温度、混合比例和来源表的条件一起报告；弛豫受分子转动相关时间和Larmor频率竞争，可非单调，不把“更黏一定T1更长”当普遍规则。未知样品数据及第三方黏度表没有冒造。

\originalsection{exp13：磁矩、化学位移与自然宽}
先用已知吸收体标速度，按相同通道的往返扫描比较驱动非线性；用能窗选14.4keV、测背景与活时间。六磁峰位置联合拟合 $E_{m_e,m_g}=E_0-g_e\mu_NBm_e+g_g\mu_NBm_g$，已知基态磁矩时可由两组分裂关系分离B与激发态g。不同铁化合物的局域场比较保留取向与厚度条件。

双四极峰、中心化学位移和温度位移各有不同来源；近薄吸收体线宽为源/吸收Lorentz宽的和，改变亚铁氰化物厚度后外推薄限可减饱和展宽。温度依赖二阶Doppler移为 $\Delta E/E\simeq-\langle v^2\rangle/(2c^2)$，固体振动不能总用高温经典 $3k_BT/M$。本项目没有源、吸收体或温控数据，故不报告新实测磁矩、宽度或相对论温度系数。

\Needspace{12\baselineskip}
\originalchapter{粒子、光谱与精密校准的测量链}
此章接续前九份手册，继续整理剩余采集与分析要求。结果仍以明确模型或真实仪器数据为前提，本项目的现场执行缺口保持原样记录。
\originalsection{exp14：muon飞行与停止寿命}
速度部分改变两探测器间距、测时间差并以电子学参考修正固定延迟；斜入射路长不是垂直间距。寿命部分识别停止与后续衰变双脉冲，保留时间窗、死区、背景和材料。对每个bin用积分期望 $A\tau(\e^{-t_i/\tau}-\e^{-t_{i+1}/\tau})+B\Delta t_i$ 作Poisson拟合，比低计数先减背景取对数稳健；前端分辨差区按预定仪器条件排除并记录。

原手册的三项大气存活讨论可参数化：10km光速飞行约33.4µs，若静止平均寿命约2.2µs则存活 $\e^{-33.4/2.2}\simeq2.6\times10^{-7}$，需要反推极大天空通量。真正运动寿命为 $\gamma\tau_0$；指定所需存活概率S，$\gamma=L/[v\tau_0(-\ln S)]$，再用能量损失修正生产能量。没有读出特定气球曲线便不报其“实测天空通量”。手册另给11g全程加减速的4光年往返讨论，理想单段距离d、固有加速度a满足 $t=(c/a)\sqrt{(1+ad/c^2)^2-1}$、固有时 $(c/a)\operatorname{arcosh}(1+ad/c^2)$；四个对称半程取d为2光年再相加，明确推进与掉头理想化。

\originalsection{exp15：Rutherford角接受与前向发散}
薄金箔、准直alpha源和硅探测器先测无靶束率及能量，再按角度测谱、活时间、背景和靶厚。绝对微分截面为扣背景计数除 $\Phi t n_t\Delta\Omega\epsilon$，原手册要求可在30度等较大角减少多重散射干扰。Rutherford $\sin^{-4}(\theta/2)$很陡，预测应平均真实接受窗，单用窗中心可有系统偏差。

原题讨论为何前向不见无限计数：纯无限程库仑势的总截面发散是无限远微小偏转的理想化；电子屏蔽截断远场、未散射束与仪器有限角宽无法分开极小偏角，多重散射和有限通量还改变可观测分布。不能把有限角积分发散直接说成探测器理应收到无限粒子。真实核近距离、厚靶能损和核力范围也限制大角解释。

\originalsection{exp17：三种谱的共同波长标定}
先扫已知汞线，拟合读数到波长关系、检查校准残差；再测氢、氘Balmer和钠双线，不在样品切换后任意重新定义零点。对氢的 $1/\lambda=R_H(1/4-1/n^2)$ 联合拟合R；共用校准协方差应传播到各线的比较。

同位素差来自约化质量 $R_M=R_\infty/(1+m_e/M)$。故 $R_D/R_H=(1+m_e/M_H)/(1+m_e/M_D)=\lambda_H/\lambda_D$，从多对线的联合比求 $M_D/M_H$，而非把波长比直接当约2的核质量比；小差会放大误差。钠双线给自旋轨道精细差，狭缝/仪器分辨不足时拟合两个卷积峰不能无条件识别；重复不同狭缝检查峰心/间距稳定性。

\originalsection{exp31：X射线的能标、壳层与吸收}
手册用能谱探测识别各元素特征线、连续谱、衰减，先能标和死时间检查，再按每条谱峰拟合背景与面积。对Kα、Kβ等光子频率平方根与Z比较Moseley标度，不能把逃逸峰或pileup峰当新元素线。不同目标/管压的激发阈值和线能分开记录。

Tb壳层讨论中K为1s、L1为2s、L2/3为2p的不同j；电偶极 $\Delta l=\pm1$，因此L1到K的s--s跃迁禁戒。可用允许的更高壳p到L1线、配合已知其他壳结合能间接求L1，不虚构一条允许的L1--K线。吸收部分每个能量、元素、厚度测 $I/I_0$，拟合 $\ln(I/I_0)=-\mu x$，报告 $\mu/\rho$ 和每原子截面 $\sigma=\mu A/(\rho N_A)$；多能束、吸收边和几何散射可使单斜率模型失败。

\originalsection{exp39：超导转变与持久流}
按手册不同探头区分AC互感转变、Meissner和铅筒持久流。测Pb、V、Nb的转变时同步保留温度、加热/冷却方向和磁场；正常金属电导突然变化也能改AC互感，所以必须结合磁通排斥，不能只凭一个电子学跳变断言超导。

铅筒高于Tc时加轴向场，带场冷却再撤外场，内部Hall读数检验保留磁通；不同冷却历史区分内外表面屏蔽电流与正常残场。Josephson/SQUID讨论用 $\Phi_0=h/(2e)$、$V_n=nhf/(2e)$连接磁通/电压标准，需真实器件与读数才可检验。本卷提供原理与分析条件，没有低温实测转变曲线或持久流结果。

\originalsection{exp43：用实际传递函数比较两种噪声}
Johnson部分改变电阻/温度，独立标定频率增益H，输出均方为 $\int_0^\infty4k_BTR|H(f)|^2\dd f$ 加放大器噪声，等效噪声带宽不能仅取标称截止差。电阻增大时输入电容与前置电流噪声都可能变大，简单直线区需用残差和带宽变化限定。

散粒部分改变平均光电流，扣暗背景后预期单边电流谱 $2eI$，输出噪声需同样的传递函数换算。对光电流的斜率估e，Johnson的斜率估kB；电子学本底可用截距/暗测，但信号相关寄生噪声不是常数截距。保留每组均方估计的独立时间块，避免把高频数字点当独立样本。该因子2来自单边谱约定，双边谱不能照抄。

\originalsection{exp46：21厘米测量到银河旋转}
射电望远镜先做指向与太阳/参考源扫描、校系统温度和增益，记UTC、方位/高度或银河经纬、频率轴、谱窗与积分时间。ON/OFF去连续基线、标HI峰，转换 $v_{\rm rad}\simeq c(\nu_0-\nu)/\nu_0$，再做地球/太阳到适当静止系速度修正。各方向无共同速度框架就不能比较银河旋转。

银河平面内内侧切点法有 $R=R_0|\sin l|$，第一象限理想圆轨道 $v_{\rm term}=V(R)-V_0\sin l$；速度包络而非任意最强峰给旋转曲线。非圆运动、光学厚度、气体分布、波束和局部干扰形成局限。FFT记录长T的频道间隔为1/T，矩形窗首零宽也为1/T；原手册一处称宽与T成正比属排印问题，真实窄化随记录变长。本项目没有现场天区谱或假造银河曲线。

\originalsection{exp48、49、51：精密频轴、脉冲和力的校准}
exp48同时扫无泵Doppler底、反向泵浦的窄凹陷及参考Fabry--Perot，用共同上能级线对之差求基态超精细，区分真实跃迁与两线中点的交叉峰。频轴可用多腔周期标定，原48cm研究级参考图的数值不充当自己的测量；扫向、模跳和激光宽进入误差预算。弱信号和功率展宽的平衡需多强度比较。

exp49先量H/C频率、J和接收相位，再扫脉宽以衰减正弦拟合90度长度；原始复积分不应只取绝对值抹掉负峰。量T1/T2、执行CNOT真值表和四函数Deutsch--Jozsa，矩阵演算与物理脉冲次序已在前章完整给出。可选Grover需要伪纯输入、四oracle及多次迭代强度，二维理想搜索 $\sin^2[(2k+1)\arcsin(1/2)]$ 在k为1时最大，再下降；本卷仅提供条件与理想预测，不冒称做了仪器Grover测量。

exp51先用固定珠二维扫描校QPD与台位置，浮珠每档功率至少重复三次；无驱动时间序列同时作协方差/equipartition和PSD，受迫X/Y扫描可用Stokes阻力交叉校力。高通、低通、采样混叠和近壁阻力改变两种估计，不能把它们不一致先归于Boltzmann定律错误。洋葱囊泡先低功率记直线速度，再逐步增陷阱至停止，停力 $k\Delta x$须用当时功率校准和重复事件；E. coli和DNA为需要提前制样的可选扩展，不把未取得材料写成已测。

\originalsection{把未采集项留在报告里}
最终报告应并列目标、原始可观察量、独立校准、模型、结果、统计及系统误差和适用边界。上述十九份手册各有可执行的教学链，但没有任何本项目现场数据；材料卷的真实DFT/MD/MC结果另章标软件和版本。Raman学生提案仅作独立补充，与正式手册、官方完成结果分开。不可用网页预习题、第三方仅书本题号和实体训练条件不会因正文详尽而自动变成已完成实验。
